% A.66 candidate. NOTES ONLY, UNREGISTERED; independent audits required.
\subsection{An abelian factor of order prime to the Heisenberg prime}

\begin{statement}{A.66}{Open}{Candidate: Heisenberg times prime-to-prime abelian closure}
Let \(p\) be an odd prime and let \(B\) be a finite abelian group
of order prime to \(p\). Every geometrically connected finite
cover of \(U=E\setminus\{O\}\), \(E:y^2=x^3-x\), whose geometric
Galois closure is a quotient of
\[
                         H_p\times B
\]
satisfies both assertions of Task A.5, at every coefficient
prime, every original cusp packet, and every permitted open
derived subgroup. Here \(H_p\) is the exponent-\(p\), order-\(p^3\)
Heisenberg group.

More concretely, take any intermediate \(D^\circ\to U\) of
A.65 and any elliptic isogeny \(\psi:E'\to E\) of degree \(h\)
prime to \(p\). The actual fibre product
\[
 Y=D^\circ\times_U
       \bigl(E'\setminus\psi^{-1}(O)\bigr)
                                                               \tag{A.66.1}
\]
is geometrically connected and has both assertions. Its proper
compactification \(D'\to D\) is finite étale of degree \(h\), and
\[
 \deg(Y/U)=h\deg(D^\circ/U),\quad
 g(D')=1+h(g(D)-1),\quad
 \#(D'\setminus Y)=h\#(D\setminus D^\circ).
                                                               \tag{A.66.2}
\]
In particular the noncentral order-\(p\) quotient gives degree
\(hp^2\), genus \(1+hp(p-1)/2\), and \(hp\) cusps; the full
Heisenberg case gives degree \(hp^3\), genus
\(1+hp^2(p-1)/2\), and \(hp^2\) cusps.

This is a fixed-subgroup theorem for actual composite covers,
not a transfer through a surjective pro-prime quotient. It does
not treat products of two nonabelian Sylow factors or all finite
nilpotent closures. The universal Task A.5 remains open.
\end{statement}

\begin{proof}[Candidate proof, pending independent audits]
All geometric covers, groups, isogenies and marked points can
first be defined over a finite number-field extension. The
comparison with every original packet and every permitted open
derived subgroup is the finite-extension comparison of
A.41--A.42. We establish the actual pointed stable-model
hypotheses before using it.

\paragraph{Coprime factors and the actual function fields.}
If \(S\le H_p\times B\), then
\[
                         S=S_p\times S_B.
                                                               \tag{A.66.3}
\]
Indeed, for \((a,b)\in S\), choose an integer that is one
modulo the order of \(a\) and zero modulo the order of \(b\);
raising \((a,b)\) to that integer isolates \((a,1)\in S\).
The reverse choice isolates \((1,b)\). This proves (A.66.3)
without a normality assumption on \(S\).

An abelian geometric quotient of the rank-two fundamental
group kills its cusp commutator, so its covering extends
étale over \(E\). After making a point above \(O\) rational
and using it as origin, this proper cover is an elliptic
isogeny. Its deck group is its kernel. The intrinsic
Heisenberg comparison in A.65 identifies the \(H_p\)-quotient
with the actual cover of A.61--A.62.

The two Galois function fields have coprime degrees. Their
intersection over the base function field is trivial and,
since they are Galois, they are linearly disjoint. Their
compositum therefore has Galois group \(H_p\times B\).
Taking fixed fields under (A.66.3) identifies each intermediate
with the actual fibre product (A.66.1), with the abelian
isogeny replaced by its quotient under \(S_B\).
This argument proves geometric connectedness as well as the
identification of the covering map. It is not an identification
made from group orders alone.

The same reasoning applies if the closure is only a quotient
of \(H_p\times B\): its normal kernel splits into the two
coprime factors. A proper nonabelian quotient of \(H_p\) does
not exist, by A.65; the remaining case is abelian and hence
has genus-one compactification. That case is already A.40.

The projection of (A.66.1) to \(D^\circ\) extends to
\[
                         D'=D\times_E E'\longrightarrow D.
                                                               \tag{A.66.4}
\]
It is finite étale because \(\psi\) is. Thus \(D'\) is smooth
proper and geometrically connected, and there are exactly
\(h\) points above every old cusp. Riemann--Hurwitz proves
all formulas (A.66.2). When \(D\) has genus one this again
falls under A.40. It remains to handle \(S_p=1\) and a
noncentral order-\(p\) subgroup.

\paragraph{The one nonconstant component in the actual models.}
In these two cases A.61--A.65 supply a factorization
\[
       D\xrightarrow{q_D} A\xrightarrow{\lambda}E,
       \qquad f_D=\lambda q_D ,
                                                               \tag{A.66.5}
\]
where \(A\) is an elliptic curve isogenous to \(E\) and
\(\lambda\) has \(p\)-power degree. For the full Heisenberg
cover, \(A=E,\ q_D=q_p,\ \lambda=[p]\). For a noncentral
subgroup, \(T\subset E[p]\) is its translation image,
\(A=E/T\), and \(\lambda=\widehat\phi\), with
\(\phi:E\to E/T\) as in (A.65.3).

Choose the place over \(p\) used in A.61 or A.62. In the
ordinary noncentral case choose the Gaussian place where
\(T\) has nonzero reduction, exactly as in A.65. After a
finite extension, \(A\) has good reduction and the pointed
semistable model \(\mathcal D\) has the following properties:

\begin{enumerate}
\item its dual graph is a tree;
\item the actual map \(q_D\) extends to the good model
\(\mathcal A\);
\item there is exactly one component \(M\) on which the
special map is nonconstant, and
\(M\to\overline A\) is finite radicial;
\item every other component maps constantly to
\(\overline A\);
\item every cusp has a distinct cusp joined to it by a path
consisting entirely of rational components, including the
components containing the two marks.
\end{enumerate}

For the full cover these are the explicit main, tail and
marking conclusions of A.61--A.62. For an inert-prime
noncentral quotient the actual resolution and domination
argument in A.65 proves (ii)--(iv), and its rational fibre
over the origin proves (v). Compact type gives (i).

At the ordinary noncentral place, property (i) follows
as well from the actual surjective norm \(J_C\to J_D\):
the source has good reduction by A.62, so the target
does too by the same Serre--Tate Corollary 2. The
semistable target graph therefore has cycle rank zero.

Here is the additional verification of (iii) for the
ordinary noncentral quotient. In A.65 the order-\(p\)
group acts freely on the entire semistable model because
its translation reduction is nonzero. It preserves the
unique nonconstant component of the source, which is an
elliptic curve \(M_C\) with radicial degree-\(p\) map
\(M_C\to\overline E\). Its quotient
\(M_C\to M_C/S_p\) is finite étale of degree \(p\).
The special isogeny
\(\overline\phi:\overline E\to\overline A\) is also
étale of degree \(p\). The commuting square
\[
 \begin{array}{ccc}
 M_C&\longrightarrow&\overline E\\
 \downarrow&&\downarrow\overline\phi\\
 M_C/S_p&\longrightarrow&\overline A
 \end{array}
                                                               \tag{A.66.6}
\]
therefore makes the bottom map degree \(p\) with separable
degree one: the composite has total degree \(p^2\) and
separable degree \(p\), and the left map is étale of degree
\(p\). Hence the bottom map is radicial. Every other source
component is constant over \(\overline A\), so no other
nonconstant target component appears. The quotient model
and its marking comparison are exactly those of A.65.
This proves all five properties in this case too.

If point blowups are used to separate marks or resolve a
map, every new component is rational and maps constantly
to the elliptic special fibre. Thus the same properties
hold on one common pointed semistable model.

\paragraph{An actual connected elliptic base change.}
Form the proper fibre product
\[
                         A'=A\times_E E',
                                                               \tag{A.66.7}
\]
using \(\lambda\) and \(\psi\). The two function fields
have coprime degrees, so \(A'\) is geometrically connected.
The projection \(A'\to A\) is finite étale of degree \(h\).
The point \((O,O)\) is an origin, making \(A'\) an elliptic
curve and this projection an isogeny.
Associativity of the actual fibre products gives
\[
                         D'=D\times_A A'.
                                                               \tag{A.66.8}
\]

At the chosen place, \(A'\) has good reduction because it is
isogenous to the good elliptic curve \(A\);
\cite[Theorem 1 and Corollary 2, p.~493]{fghSerreTate1968}.
The isogeny and its dual extend between the good abelian
schemes. Their composite is multiplication by \(h\), which
is étale since \(p\nmid h\). Invariant differentials show
that both extended maps have invertible differential.
Their kernels are finite, being contained in the
finite étale \(h\)-torsion. Properness makes the maps
finite, and they are étale. Their degree is the generic
degree \(h\), by flatness. Thus the actual projection
extends as a finite étale isogeny
\[
                    \mathcal A'\longrightarrow\mathcal A .
                                                               \tag{A.66.9}
\]
This is a comparison of good elliptic models, not a
prime-to-prime assertion about an arbitrary singular source.

\paragraph{The pulled-back graph and the marked paths.}
Take the actual model
\[
              \mathcal D'=\mathcal D
                              \times_{\mathcal A}\mathcal A'.
                                                               \tag{A.66.10}
\]
It is proper and semistable because it is finite étale
over \(\mathcal D\). Its generic fibre is the connected
curve (A.66.8). We work over an algebraically closed
residue field to describe its components.

The component over \(M\) is connected. Indeed
\[
                   M'=M\times_{\overline A}\overline A'
\]
maps radicially and surjectively to the connected elliptic
curve \(\overline A'\). A finite radicial surjection is a
universal homeomorphism, and this statement is preserved
by base change. The map \(M'\to M\) is étale; hence
\(M'\) is smooth. There is exactly one component over \(M\).

Every other component of \(\mathcal D\) maps to a single
point of \(\overline A\), so its pullback is a disjoint
union of \(h\) identical components. Removing the vertex
of \(M\) from the tree graph of \(\mathcal D\) leaves
subtrees, each attached to \(M\) in one node. The map to
\(\overline A\) is constant on each whole subtree: adjacent
constant components agree at their common node.
Consequently its pullback consists of \(h\) disjoint
copies, attached at the \(h\) points above that original
node on \(M'\). It follows that the whole special graph
of \(\mathcal D'\) is a tree. In particular the new
Jacobian has good reduction; no cycle is created by the
prime-to-\(p\) base change.

Now fix any marked cusp of \(\mathcal D'\). The chosen
old cusp lies in a rational joining path with a distinct
old cusp, by property (v). The elliptic map is constant
on this entire path. Its pullback is therefore \(h\)
disjoint identical copies. The selected new cusp belongs
to one copy, and the other endpoint of that copy is a
distinct marked cusp. Both endpoints and the entire
joining path lie on rational components. This argument
works for every lift of every cusp, not just for one
chosen component.

Equivalently, any subsequent stabilization contracts only
rational components, and marked sections can be separated
on the same rational trees. The actual pointed criterion
of A.41--A.42 is therefore satisfied by \(\mathcal D'\).
It proves the full fixed-subgroup equality at every
coefficient prime, including the residue prime \(p\),
and for every permitted packet and open derived subgroup.

\paragraph{Normalizers and scope.}
The geometric closure of (A.66.1) is a subgroup quotient
of \(H_p\times B\), hence is nilpotent. A.59 supplies the
normalizer assertion at every coefficient prime. The
fixed-subgroup argument above used the actual model,
not this nilpotence alone. In particular it supplies no
claim about a product of two nonabelian prime factors,
where the additional cover is not a proper elliptic
isogeny and the constant-component pullback argument
does not apply.
\end{proof}
